\subsection{其他连接词}
\label{other_connectives}

为了使表达式更加简洁，再定义一些“派生连接词”。需要注意的是，这些连接词都不是 $\mathcal{L}_0$ 的初始连接词，而是通过 $\neg$ 和 $\to$ 定义出来的，因此，它们并不增加 $\mathcal{L}_0$ 的表达能力。

\begin{definition}[合取 Conjunction]
    \label{conjunction_of_L0}
    \index{合取 Conjunction}
    设 $\varphi$ 和 $\psi$ 是 $\mathcal{L}_0$ 的公式。$\varphi$ 和 $\psi$ 的合取记为 $\varphi \wedge \psi$，定义为：
    \[
        \varphi \wedge \psi := \neg (\varphi \to \neg \psi)
    \]
\end{definition}

\begin{definition}[析取 Disjunction]
    \label{disjunction_of_L0}
    \index{析取 Disjunction}
    设 $\varphi$ 和 $\psi$ 是 $\mathcal{L}_0$ 的公式。$\varphi$ 和 $\psi$ 的析取记为 $\varphi \vee \psi$，定义为：
    \[
        \varphi \vee \psi := (\neg \varphi \to \psi)
    \]
\end{definition}

\begin{definition}[双条件 Biconditional]
    \label{biconditional_of_L0}
    \index{双条件 Biconditional}
    设 $\varphi$ 和 $\psi$ 是 $\mathcal{L}_0$ 的公式。$\varphi$ 和 $\psi$ 的双条件记为 $\varphi \leftrightarrow \psi$，定义为：
    \[
        \varphi \leftrightarrow \psi := ((\varphi \to \psi) \wedge (\psi \to \varphi))
    \]
\end{definition}

\vspace{1em}
\textbf{括号的省略约定}：为了使表达式更加简洁，在不引起歧义的前提下，引入括号的省略约定：
\begin{enumerate}
    \item 公式最外层的括号可以省略
    \item 连接词的结合力依次递减为： $\neg$，$\wedge$，$\vee$，$\to$，$\leftrightarrow$
    \item 相同优先级的连接词，从右向左结合
    \item 其他情况下，括号不能省略
\end{enumerate}
\vspace{1em}

\begin{example}
    以下是一些根据上述省略约定得到的公式：
    \begin{enumerate}
        \item - $p\to q \to r$ 等价于 $p\to (q \to r)$
        \item - $p\wedge q \vee r$ 等价于 $(p\wedge q) \vee r$
        \item - $r\to p \to p \leftrightarrow \neg \neg p \vee q$ 等价于 $(r\to (p \to p)) \leftrightarrow (\neg \neg p \vee q)$
    \end{enumerate}
\end{example}